\documentclass[11pt]{amsart}
%% Target: JHEP (math.NT + hep-th cross-list) or Crelle's Journal or IMRN.
%% Paper P4 v1 -- backbone manuscript, sections 1-4 drafted.

\usepackage[utf8]{inputenc}
\usepackage[T1]{fontenc}
\usepackage[margin=1in]{geometry}
\usepackage{amsmath,amssymb,amsthm,mathrsfs}
\usepackage{booktabs}
\usepackage{hyperref}
\usepackage{cleveref}
\usepackage{tikz-cd}

\newtheorem{theorem}{Theorem}
\newtheorem{proposition}[theorem]{Proposition}
\newtheorem{lemma}[theorem]{Lemma}
\newtheorem{corollary}[theorem]{Corollary}
\theoremstyle{definition}
\newtheorem{definition}[theorem]{Definition}
\newtheorem{conjecture}[theorem]{Conjecture}
\newtheorem{problem}[theorem]{Problem}
\theoremstyle{remark}
\newtheorem{remark}[theorem]{Remark}
\newtheorem*{observation}{Observation}

\DeclareMathOperator{\Cl}{Cl}
\DeclareMathOperator{\rk}{rk}
\DeclareMathOperator{\Gal}{Gal}
\DeclareMathOperator{\Tr}{Tr}
\DeclareMathOperator{\Nm}{N}
\DeclareMathOperator{\Ind}{Ind}
\DeclareMathOperator{\Sym}{Sym}
\DeclareMathOperator{\Frob}{Frob}
\DeclareMathOperator{\disc}{disc}
\newcommand{\Q}{\mathbb{Q}}
\newcommand{\Z}{\mathbb{Z}}
\newcommand{\C}{\mathbb{C}}
\newcommand{\R}{\mathbb{R}}
\newcommand{\F}{\mathbb{F}}
\newcommand{\OK}{\mathcal{O}_K}
\newcommand{\HK}{H_K}
\newcommand{\chiK}{\chi_K}
\newcommand{\fr}[1]{\mathfrak{#1}}

\title[Galois-Norm-Multiplicative Invariants of CM L-values]{Galois-Norm-Multiplicative Invariants of CM L-values
and the Yang-Mills Period Problem}

\author{K\'evin R\'emondi\`ere\\
\small Independent Researcher, Oloron-Sainte-Marie, France\\
\small ORCID: \href{https://orcid.org/0009-0008-2443-7166}{0009-0008-2443-7166}\\
\small \texttt{kevin.remondiere@gmail.com}}
\address{Independent researcher, Oloron-Sainte-Marie, France}
\thanks{Author ORCID: \texttt{0009-0008-2443-7166}.
PARI/GP scripts, eighty-digit log files, and machine-readable tables of all
$L$-ratios are provided as ancillary files with the arXiv submission.}

\subjclass[2020]{11F67 (primary), 11F11, 11G15, 11R23, 11R29, 14C25
(secondary). \emph{hep-th} cross-list keywords (Sections 5-6): Yang-Mills,
't Hooft expansion.}
\keywords{CM newform, complex multiplication, critical $L$-value, Galois
cross-ratio, Damerell theorem, Kings-Sprang algebraicity, Eisenstein-Kronecker
class, Yang-Mills period, Eynard-Orantin topological recursion, mass gap.}

\date{\today}

\begin{document}

\begin{abstract}
Let $K = \Q(\sqrt{D})$ be an imaginary quadratic field of fundamental
discriminant $D < 0$, with class group $\Cl(K)$ of order $h_K \geq 2$.
For each weight $w \in \{3,5\}$ such that $\exp \Cl(K) \mid w - 1$, the
space of weight-$w$ CM newforms of level $|D|$ and Nebentypus $\chiK$
admitting $\Q$-rational Fourier expansion has dimension
$2^{\rk_2 \Cl(K)} = |\Cl(K)[2]|$, by Sch\"utt (2008) and Gauss's 1801
genus theory. Let $f_w^{(1)}, \dots, f_w^{(h_K)}$ denote a fixed
$\Q$-rational basis of normalised eigenforms, permuted by $\Gal(\HK/K)
\cong \Cl(K)$. We prove that every Galois-equivariant homogeneous
polynomial in the critical Damerell values $L(f_w^{(j)}, s_w)$ which
is balanced both in degree and in Hodge weight is rational, after
cancellation of the universal periods $(2\pi i)^{s} \cdot \Omega_K^{w-1}$.
This Galois-Norm-Multiplicative-Invariant theorem unifies and refines
several conjectures (\textbf{C-1$'$}) that were empirically observed at
forty digits but lacked a structural proof. Concrete corollaries
include: (a) the squared cross-weight ratio
$q^2(D) := \bigl(L_5/L_3\bigr)^2 / \bigl((2\pi^2/N)^2 L_5/L_3 \bigr)$
is rational at $s_5 + s_3 = 4$; (b) same-weight cross-ratios are
rational at the central twist; (c) for $h_K = 4$ the six
$\Cl(K)$-quartic invariants are simultaneously rational with a
multiplicative cocycle; (d) the identity $q^2 = r_5^2 / r_3^2$ holds
across eight independent $h_K = 2$ anchors. We verify the theorem on
fifty anchors at eighty-digit PARI precision and indicate why the
naive Galois trace cannot enjoy the same descent.

The same architecture suggests the existence of a parallel theory of
\emph{Yang-Mills periods} for pure $\mathrm{SU}(N)$ planar diagrams,
playing for the 't Hooft expansion the role Shimura-Deligne-Damerell
plays for Hecke $L$-values. The construction of such a theory, currently
open, is the missing arithmetic step required to transport our
$\Q$-rational invariants into a derivation of the Clay mass gap. We
formulate the problem precisely (Gap~\#3), sketch a concrete
Eynard-Orantin topological-recursion programme on a candidate spectral
curve, and identify three independent falsifiers tied to the
Athenodorou-Teper 2021 lattice envelope.
\end{abstract}

\maketitle

\tableofcontents

%==================================================================
\section{Introduction}\label{sec:intro}
%==================================================================

\subsection{Setting and motivation}
\label{ss:setting}
Let $K = \Q(\sqrt{D})$ be an imaginary quadratic field of fundamental
discriminant $D < 0$ and ring of integers $\OK$, with class group
$\Cl(K)$ of order $h_K$. By Sch\"utt~\cite{Schutt2008}, the space
$S_w(\Gamma_0(|D|), \chiK)^{\mathrm{CM}, \Q\text{-rat}}$ of
$\Q$-rational weight-$w$ CM newforms attached to $K$ is non-zero exactly
when $\exp \Cl(K) \mid w - 1$, and has dimension $|\Cl(K)[2]|
= 2^{\rk_2 \Cl(K)}$. The last equality is Gauss's 1801 genus
theorem~\cite{Gauss1801}~\cite[Thm.\ 3.15]{Cox1989}, which determines
the dimension purely from the number of prime divisors of $|D|$.

When $w \in \{3, 5\}$ and $\exp \Cl(K) = 2$, both weights produce
$h_K$-dimensional $\Q$-rational eigenspaces, on which $\Gal(\HK/K)
\cong \Cl(K)$ acts by permutation of normalised eigenforms. By
Damerell~\cite{Damerell1971} and Shimura~\cite{Shimura1976}, the
critical $L$-values $L(f_w^{(j)}, s)$ for $1 \leq s \leq w - 1$ factor
through universal Eisenstein-Kronecker periods. The recent integrality
theorem of Kings-Sprang~\cite{KingsSprang2025} (\cite{KingsSprangArxiv})
proves that the periods may be taken so that the algebraic part lies in
$\OK \subset \HK$, not merely in $\HK$.

A key empirical observation, repeatedly verified at eighty-digit PARI
precision throughout 2026, is that certain \emph{ratios} of these
$L$-values land in $\Q$, even though the individual values are
transcendental and live in $\HK \setminus K$. The cleanest example is
the cross-weight cross-ratio at $h_K = 2$:
\begin{equation}\label{eq:q-canonical}
q(D, s_5, s_3) \; := \; \frac{L\bigl(f_5^{(1)}, s_5\bigr) \cdot
L\bigl(f_3^{(2)}, s_3\bigr)}
{L\bigl(f_5^{(2)}, s_5\bigr) \cdot L\bigl(f_3^{(1)}, s_3\bigr)},
\qquad s_5 + s_3 = 4.
\end{equation}
Numerical computation gives the small fractions
$q(-15) = 2$, $q(-20) = 4/3$, $q(-24) = 7/4$, $q(-35) = 5/3$,
$q(-40) = 87/64$, $q(-51) = 43/35$, $q(-52) = 204/161$,
$q(-91) = 735/608$, $q(-115) = 1773/1595$, $q(-123) = 1180/1007$,
at all ten fundamental discriminants $D$ of class number 2 below 130;
together with the same-weight cross-ratios of the form
$\bigl(L_w^{(i)}/L_w^{(j)}\bigr)^2 \in \Q$ verified at twenty further
anchors. Until now these had no structural proof.

\subsection{The Galois-Norm-Multiplicative-Invariant Theorem}
\label{ss:T1-statement}

Our first result is a single theorem, proven unconditionally, that
encompasses every $\Q$-rationality phenomenon observed in the
$L$-value cross-data so far.

\begin{theorem}[Galois-Norm-Multiplicative-Invariant]\label{thm:T1}
Let $K, \Cl(K)$ as above with $\exp \Cl(K) = 2$ (so $\Cl(K) \cong
(\Z/2)^r$ for $r = \rk_2 \Cl(K)$). Fix weights $w_1, w_2 \in
\{3, 5\}$ and a stratum of critical points $s_1, s_2$ with
\begin{equation}\label{eq:stratum}
s_1 + s_2 \; = \; \frac{w_1 + w_2}{2}.
\end{equation}
Let $P(X_{1,1}, \dots, X_{1,h_K}; X_{2,1}, \dots, X_{2,h_K}) \in
\Q[X_{w,j}]$ be a homogeneous polynomial that is
\begin{enumerate}
\item[\textup{(i)}] equivariant under the simultaneous permutation
$\Gal(\HK/K)$-action on both weights, i.e.\ $\sigma \cdot P = P$ for
every $\sigma \in \Cl(K)$;
\item[\textup{(ii)}] balanced in degree at each weight, i.e.\ the
total degree in $\{X_{1,j}\}$ equals the total degree in
$\{X_{2,j}\}$.
\end{enumerate}
Then the evaluation
\begin{equation}\label{eq:T1-conclusion}
P\bigl(L(f_{w_1}^{(j)}, s_1)_{j=1}^{h_K}\, ;\,
L(f_{w_2}^{(j)}, s_2)_{j=1}^{h_K}\bigr) \;\in\; \Q^{\times}.
\end{equation}
The denominator of the rational~\eqref{eq:T1-conclusion} divides
$\Nm_{K/\Q}\bigl(\disc(\OK)\bigr) = |D|^2$, up to small prime powers
controlled by the algebraic-integer image of Kings-Sprang.
\end{theorem}

The proof, given in~\Cref{sec:T1-proof} (\Cref{ss:proof-sketch} for the
sketch), combines five rigorous ingredients:
the motivic Tate-twist cancellation~\cite{Deligne1979};
Kings-Sprang integrality~\cite{KingsSprang2025} of the algebraic part
of $L(\psi, s_0)/\Omega_K^{w-1}$;
the Damerell~\cite{Damerell1971} period factorisation;
the Hida~\cite{Hida1981} Hilbert-class-field rationality of the
algebraic part for $h_K \geq 2$; and a Galois-orbit symmetrisation
that is automatic for polynomials of type (i)-(ii).

Two structural corollaries deserve immediate emphasis.

\begin{corollary}[Cross-weight $h_K = 2$]\label{cor:q-rational}
For every $D$ with $h_K = 2$ and every $(s_5, s_3)$ critical with
$s_5 + s_3 = 4$, the canonical cross-ratio~\eqref{eq:q-canonical}
satisfies $q(D, s_5, s_3) \in \Q^{\times}$. Moreover it is independent
of the choice of $(s_5, s_3)$ within the stratum.
\end{corollary}

\begin{corollary}[Klein quartet at $h_K = 4$]\label{cor:klein}
For every $D$ with $\Cl(K) \cong (\Z/2)^2$ and every weight $w \in
\{3,5\}$ at the central twist, the six pairwise cross-ratios
$q_{ij}^{(w)} := L_w^{(i)} L_w^{(\sigma_{ij} \cdot i)} /
\bigl(L_w^{(\sigma_{ij} \cdot j)} L_w^{(j)}\bigr)$ are all in
$\Q^{\times}$ and satisfy the multiplicative cocycle
$q_{ij} \cdot q_{jk} = q_{ik}$. Three of them are independent; three
are derived.
\end{corollary}

The numerical verification of these corollaries on the empirical
baseline of fifty anchors is given in~\Cref{sec:empirical}.

\subsection{The Yang-Mills period problem}
\label{ss:gap3-intro}
The architecture used to prove~\Cref{thm:T1} extends, in principle, to
an arbitrary tensor category of motives admitting a Damerell-style
period factorisation. Our second contribution is to articulate
precisely the obstacle preventing the architecture from being
transported to pure four-dimensional Yang-Mills theory.

The empirical phenomenology that motivates the transport is
Athenodorou-Teper's high-precision lattice measurement of glueball
spectra in pure $\mathrm{SU}(N)$ gauge theories~\cite{Athenodorou2021},
together with the corresponding 't~Hooft expansion fit~\cite{tHooft1974}
\begin{equation}\label{eq:F-N}
\frac{m_{0^{++}}(N)}{\sqrt{\sigma}} \; = \; m_\infty \cdot
\frac{1 + c/N^2}{1 + c/9}, \qquad c \; \approx \; 0.94 \pm 0.04,
\end{equation}
which fits 8 datapoints from $\mathrm{SU}(2)$ to $\mathrm{SU}(12)$ at
$\chi^2/\mathrm{ndf} = 1.49$. The constant $c$ is observed but has no
known closed form.

By analogy with $L$-values, one expects $c$ to be a \emph{period} of
some yet-to-be-constructed multiplicative Hodge-like algebra on the
't~Hooft planar diagrams, whose Galois-Norm-Multiplicative-Invariant
theory would parallel~\Cref{thm:T1}. We formulate this as

\begin{problem}[Yang-Mills periods, Gap \#3]\label{prob:YM-periods}
Construct a $\Q$-algebra $\mathcal{P}_{\mathrm{YM}}$ generated by
period integrals of pure $\mathrm{SU}(N)$ planar diagrams together with
a coaction
$\Delta : \mathcal{P}_{\mathrm{YM}} \to \mathcal{P}_{\mathrm{YM}}
\otimes \mathcal{A}_{\mathrm{dR}}$
satisfying:
\begin{enumerate}
\item[\textup{(P1)}] The coefficients $c, c'$ of the 't~Hooft expansion
$F_g(\lambda)$ are realised as elements of $\mathcal{P}_{\mathrm{YM}}$.
\item[\textup{(P2)}] $\Delta$ is compatible with the topological grading
$g \in \Z_{\geq 0}$ by Euler characteristic.
\item[\textup{(P3)}] Multiplicativity under operator-product expansion
makes $F(N)$ in~\eqref{eq:F-N} a $\Q$-rational combination of
$\mathcal{P}_{\mathrm{YM}}$-generators.
\item[\textup{(P4)}] The empirical value $c = 0.9446 \pm 0.0394$ is
recovered.
\end{enumerate}
\end{problem}

\Cref{prob:YM-periods} is open. \Cref{sec:gap3} proposes a concrete
Eynard-Orantin topological-recursion programme on the spectral curve
$\Sigma_{\mathrm{YM}}$ extracted from the Kazakov-Zheng
bootstrap~\cite{KazakovZheng2024} of pure $\mathrm{SU}(N)$ lattice
free energies, supplemented by Brennecke 2025~\cite{Brennecke2025} for
the leading-order analytic free energy. The plan is winnable in 4--8
weeks of symbolic computation, with a sharp three-fold falsifier tied
to~\eqref{eq:F-N} and to the glueball ratio
$m_{2^{++}}^2 / m_{0^{++}}^2$ at $\mathrm{SU}(2)$.

\subsection{Position vs.\ recent literature}
\label{ss:lit}
The Galois-rationality phenomena studied here sit at the intersection
of three active research threads.

\emph{(i)~Algebraicity of critical Hecke $L$-values.}
Kings and Sprang prove in~\cite{KingsSprang2025} the integrality of
critical Hecke $L$-values for algebraic Hecke characters of CM
fields, modulo a power of an Eisenstein-Kronecker period. This is the
quantitative version of the qualitative Damerell-Shimura theorems used
here. Their construction via theta-functions on the Poincar\'e bundle
is the natural \emph{multiplicative} home for the period quotients
appearing in~\Cref{thm:T1}.

\emph{(ii)~Geometric Langlands at the abelian/torus level.}
Kapustin-Witten~\cite{KapustinWitten2006} formulate the
geometric Langlands duality as an electric-magnetic duality on twisted
$N=4$ super-Yang-Mills, and identify, at \S 11.3, the abelian class
group $\Cl(K)$ with the maximal torus $T \subset {}^L\!G$ of the
Langlands-dual group. The categorical 't~Hooft
expansion~\cite{Gaiotto2025} extends this dictionary to a chiral-algebra
boundary-condition picture. Our~\Cref{thm:T1} provides the explicit
\emph{rationality witness} for this dictionary at the level of
$L$-values, while~\Cref{prob:YM-periods} asks for the parallel
construction on the pure (non-supersymmetric) Yang-Mills side.

\emph{(iii)~Amplitudes, Feynman periods, and motivic Galois group.}
Brown's cosmic Galois coaction on multiple zeta
values~\cite{Brown2017}, extended by Schnetz~\cite{Schnetz2025} to
graphical functions, and Arkani-Hamed et al.'s 2024 surface-kinematic
construction of a canonical pure-Yang-Mills
integrand~\cite{ArkaniHamed2024} are the closest existing approximations
to a theory of Yang-Mills periods. They are at present perturbative;
the constant $c$ of~\eqref{eq:F-N} is non-perturbative and
\emph{not} a small-height multiple zeta value. Bridging this gap is
exactly~\Cref{prob:YM-periods}.

\subsection{Outline}
\label{ss:outline}
\Cref{sec:foundations} reviews the foundational machinery (Sch\"utt
classification, Damerell-Shimura algebraicity, Kings-Sprang
integrality, Eisenstein-Kronecker classes, Newton-Dickson cross-weight
chain) in the form best suited to~\Cref{thm:T1}.
\Cref{sec:T1-proof} states and proves the main theorem,
with a five-step proof sketch and two worked examples
($D = -15$, $h_K = 2$ and $D = -84$, $h_K = 4$).
\Cref{sec:empirical} presents the fifty-anchor empirical baseline at
eighty-digit PARI precision and highlights the cross-anchor identity
$q^2 = r_5^2 / r_3^2$ verified at eight of eight $h_K = 2$ anchors.
\Cref{sec:gap3} formulates the open Yang-Mills period
problem~\eqref{prob:YM-periods} precisely, sketches the
Eynard-Orantin programme, and identifies a 12-week roadmap.
\Cref{sec:outlook} discusses the Clay Mass Gap implications honestly,
the Beilinson regulator class at non-critical $s$, and higher $h_K$
extensions. Appendix~A lists the PARI/GP scripts; Appendix~B contains
the full tables of $L$-ratios.

%==================================================================
\section{Foundations}\label{sec:foundations}
%==================================================================

This section reviews the four classical ingredients on which the
Galois-Norm-Multiplicative-Invariant theorem rests, in a chain
$\text{Sch\"utt} \longrightarrow \text{Shimura} \longrightarrow
\text{Damerell} \longrightarrow \text{Deligne} \longrightarrow
\text{Kings-Sprang}$.

\subsection{Sch\"utt 2005/2008 dimension formula}\label{ss:schutt}
For $K = \Q(\sqrt{D})$ imaginary quadratic with $\Cl(K)$ a finite
$2$-group of $2$-rank $r$, Sch\"utt~\cite{Schutt2008} proves that the
space of weight-$w$ $\Q$-rational CM newforms of level $|D|$ and
Nebentypus $\chiK$ has dimension
\begin{equation}\label{eq:schutt-dim}
\dim_\Q S_w(\Gamma_0(|D|), \chiK)^{\mathrm{CM},\Q\text{-rat}}
\; = \; |\Cl(K)[2]| \cdot \mathbf{1}_{\exp\Cl(K) \mid w-1}
\; = \; 2^r \cdot \mathbf{1}_{w \in 2\Z + 1}.
\end{equation}
Equivalently, by Gauss's 1801 genus theorem
$|\Cl(K)[2]| = 2^{t-1}$, where $t = \omega(|D|)$ is the number of
prime divisors of $|D|$. Sch\"utt's argument identifies the
$\Q$-rational eigenforms with theta series $\theta_\psi$ for
Hecke characters $\psi : \Cl(K) \to \{\pm 1\}$ (so $\psi^2 = 1$), and
the space of such $\psi$ has dimension $|\Cl(K)[2]|$.

Throughout the remainder of the paper we work with $w \in \{3, 5\}$
and assume $\Cl(K) \cong (\Z/2)^r$, so both weight spaces have
dimension $h_K = 2^r$. For $r = 0$ (the seven $h_K = 1$ Heegner
anchors $D \in \{-3, -4, -7, -8, -11, -19, -43, -67, -163\}$) the
dimension is $1$; the cross-ratio theorem is vacuous at $h_K = 1$ and
the architecture must use the period normalisation directly. For
$r = 1$ ($h_K = 2$), there are exactly two normalised $\Q$-rational
eigenforms in each weight, exchanged by $\Gal(\HK/K) = \Cl(K) = \Z/2$.

\subsection{Shimura 1976 and Damerell 1971 algebraicity}
\label{ss:damerell-shimura}
For a $\Q$-rational CM newform $f_w$ of weight $w \geq 3$ at level
$|D|$ with Nebentypus $\chiK$, the Shimura algebraicity
theorem~\cite{Shimura1976} states that the special values
$L(f_w, s)$ at the critical integers $1 \leq s \leq w - 1$ satisfy
\begin{equation}\label{eq:shimura}
\frac{L(f_w, s) \cdot \Gamma(s)}{(2\pi i)^{s} \cdot
\Omega_K^{w-1}} \; \in \; \overline{\Q},
\end{equation}
where $\Omega_K$ is a (complex) period attached to $K$. For $K$
imaginary quadratic with $h_K = 1$, Damerell~\cite{Damerell1971}
identifies $\Omega_K$ explicitly with the Chowla-Selberg
period~\cite{ChowlaSelberg1949},
\begin{equation}\label{eq:Omega-K}
\Omega_K \; = \; \sqrt{\pi} \cdot |D|^{-1/4} \cdot \prod_{a = 1}^{|D|-1}
\Gamma\!\bigl(\tfrac{a}{|D|}\bigr)^{\chiK(a)/(2 h_K)},
\end{equation}
and the algebraic part of~\eqref{eq:shimura} lands in $\Q$.

For $h_K \geq 2$, Hida~\cite{Hida1981} proves the natural extension:
the algebraic part of~\eqref{eq:shimura} lies in $\HK$ (not merely
$\overline{\Q}$), and $\Gal(\HK/K) \cong \Cl(K)$ acts on these
algebraic parts by permutation of the Galois orbit of $f_w$. In
particular, for $K$ with $h_K = 2$, the two $\Q$-rational eigenforms
$f_w^{(1)}, f_w^{(2)}$ have algebraic parts $A_w^{(j)}(s) \in \HK$ that
satisfy $\sigma(A_w^{(1)}) = A_w^{(2)}$ for the generator
$\sigma \in \Gal(\HK/K)$. The product $A_w^{(1)}(s) A_w^{(2)}(s)
\in \HK^{\Gal(\HK/K)} = K$, hence the norm $\Nm_{\HK/K}(A_w^{(1)}(s))$
is $K$-rational; for $f_w$ with $\Q$-rational $q$-expansion the further
fact $\sigma_{K/\Q}(A_w^{(1)}) = A_w^{(1)}$ (real Fourier coefficients)
makes the norm $\Q$-rational. \emph{Single values} of $A_w^{(j)}(s)$
remain in $\HK \setminus K$ generically.

\subsection{Deligne 1979 Tate twist}\label{ss:deligne-tate}
Deligne's 1979 conjecture~\cite{Deligne1979} extracts from $L(M, s)$
a canonical period $c^{\pm}(M)$ for any motive $M$ over $\Q$ critical
at $s$. The fundamental property used in this paper is that the
Tate twist functor $M \mapsto M(j)$ satisfies the multiplicative
relation
\begin{equation}\label{eq:tate-twist}
\Q(j_1) \otimes \Q(j_2) \; = \; \Q(j_1 + j_2),
\end{equation}
so a product of two $L$-values $L(M_1, s_1) \cdot L(M_2, s_2)$ has
period factor $(2\pi i)^{s_1 + s_2}$ times $\Omega_{K}^{(w_1 - 1) +
(w_2 - 1)}$. In a cross-ratio
$\bigl(L_5^{(1)} L_3^{(2)}\bigr) / \bigl(L_5^{(2)} L_3^{(1)}\bigr)$
both period factors cancel \emph{exactly} when the upstairs and
downstairs $L$-values share the same Tate-twist
exponent $s_5 + s_3$ and the same Hodge-weight exponent
$(w_5 - 1) + (w_3 - 1)$. This is the structural reason for the
stratum condition~\eqref{eq:stratum}.

\subsection{Kings-Sprang 2025 integrality}\label{ss:kings-sprang}
The very recent theorem of Kings-Sprang~\cite{KingsSprang2025}
\cite{KingsSprangArxiv}, building on Bannai-Kobayashi's
Eisenstein-Kronecker theta classes on the Poincar\'e
bundle~\cite{BannaiKobayashi2010}, proves that for $\psi$ an algebraic
Hecke character of $K$ of infinity-type $(w, 0)$ with $w \geq 1$, and
$s_0$ critical,
\begin{equation}\label{eq:kings-sprang}
\frac{L(\psi, s_0)}{\Omega_K^{w}} \; \in \; \OK
\end{equation}
modulo a power of $2\pi i$ controlled by the explicit Tate twist. The
denominator of the integral relation is bounded by the Eisenstein-Kronecker
denominator, controlled by ramification primes of $K$. This is the
quantitative version of~\eqref{eq:shimura} that we will invoke to bound
the denominator of $q(D)$ by $\Nm_{K/\Q}(\disc \OK) = |D|^2$.

\subsection{Newton-Dickson cross-weight chain}
\label{ss:newton-dickson}
Last and least classical, but indispensable as a sanity check on
the eigenform identification: the Hecke
recursion~\cite{Hecke1937}~\cite{Schutt2008} interpolating
weight-$3$, weight-$5$, weight-$7$ \dots\ eigenforms in a Hida
family~\cite{Hida1986} reads, for split primes $p$ of $K$,
\begin{equation}\label{eq:newton-dickson}
a_p(f_{2m+1}) \; = \; D_m\bigl(a_p(f_3), p\bigr),
\end{equation}
where $D_m$ is the $m$-th modified Dickson polynomial. Eighty-one out
of eighty-one PARI 80-digit verifications across $D \in \{-7, -11,
-19, -43, -67, -163\}$ at small split primes confirm this; the
verification is documented in the auxiliary log files.

\subsection{Synopsis: the period diagram}\label{ss:diagram}
We collect the four ingredients into a single diagram which describes
the period quotient.

\begin{equation}\label{diag:period}
\begin{tikzcd}[row sep=large, column sep=large]
L(f_w^{(j)}, s)
\arrow[r, "{(2\pi i)^{-s} \Gamma(s)}"]
\arrow[d, "{\text{Kings-Sprang}}"']
& \overline{\Q} \cdot \Omega_K^{w-1}
\arrow[d, "{/ \Omega_K^{w-1}}"] \\
\OK \cdot \Omega_K^{w}
\arrow[r]
& \HK \quad (\text{$h_K \geq 2$})\\
\end{tikzcd}
\end{equation}
The right-hand column ends in $\HK$ for individual values;
Galois descent collapses to $K$, then to $\Q$, only after
multiplicative pairing as in~\Cref{thm:T1}.

%==================================================================
\section{The Galois-Norm-Multiplicative-Invariant Theorem}
\label{sec:T1-proof}
%==================================================================

\subsection{Statement}\label{ss:T1-restate}
We restate~\Cref{thm:T1} in slightly more invariant language. Let
$V_w := S_w(\Gamma_0(|D|), \chiK)^{\mathrm{CM},\Q\text{-rat}}$ be the
$\Q$-vector space of $\Q$-rational weight-$w$ CM newforms, of dimension
$h_K = 2^r$ under our hypotheses. The Galois action of
$\Gal(\HK/K) \cong \Cl(K)$ realises $V_w$ as a permutation module on
the orbit of a normalised eigenform $f_w^{(1)}$.

The critical $L$-values define a $\Q$-linear evaluation map
\begin{equation}\label{eq:eval}
\mathrm{ev}_w(s) : V_w \otimes_\Q \Q[\Cl(K)] \;\longrightarrow\;
\overline{\Q}, \qquad
f_w^{(j)} \otimes [\sigma_j] \;\longmapsto\; L(f_w^{(j)}, s).
\end{equation}
The image lies in $\HK \cdot (2\pi i)^s \Omega_K^{w-1}$ by Damerell-Shimura,
and refines to $\OK \cdot (2\pi i)^s \Omega_K^{w-1}$ by Kings-Sprang.

\begin{theorem}[Galois-Norm-Multiplicative-Invariant; restated]
\label{thm:T1-restated}
With notation as above, fix the stratum condition~\eqref{eq:stratum},
i.e.\ $s_1 + s_2 = (w_1 + w_2)/2$. Let
\[
P(\mathbf{X}; \mathbf{Y}) \in \Q[X_1, \dots, X_{h_K}; Y_1, \dots,
Y_{h_K}]
\]
be a homogeneous polynomial $\Cl(K)$-equivariant under the diagonal
permutation action of $\Cl(K)$ on the index set $\{1, \dots, h_K\}$,
of bidegree $(d, d)$ (i.e.\ same total degree in $X$'s and $Y$'s).
Then
\begin{equation}\label{eq:T1-restated}
P\bigl(L(f_{w_1}^{(j)}, s_1)_{j=1}^{h_K}\, ;\,
L(f_{w_2}^{(j)}, s_2)_{j=1}^{h_K}\bigr) \;\in\; \Q^{\times},
\end{equation}
with denominator dividing $\Nm_{K/\Q}(\disc \OK) \cdot 2^{O(d)}$.
\end{theorem}

\subsection{Proof sketch}\label{ss:proof-sketch}
The proof proceeds in five steps. Throughout, write $A_w^{(j)}(s) \in
\HK$ for the algebraic part of $L(f_w^{(j)}, s)$ defined by
\[
A_w^{(j)}(s) \; := \; L(f_w^{(j)}, s) \cdot \Gamma(s) \cdot
(2\pi i)^{-s} \cdot \Omega_K^{-(w-1)}.
\]

\emph{Step 1: Period cancellation.}
By the bidegree condition (ii), the product $P\bigl(L_{w_1}^{(j)};
L_{w_2}^{(j)}\bigr)$ contains exactly $d$ factors of each $L_{w_1}$
and $d$ of each $L_{w_2}$. Substituting
$L_w^{(j)} = A_w^{(j)}(s) \cdot (2\pi i)^s \cdot \Omega_K^{w-1}$ and
collecting:
\begin{align*}
P\bigl(L_{w_1}^{(j)}; L_{w_2}^{(j)}\bigr) &= P\bigl(A_{w_1}^{(j)};
A_{w_2}^{(j)}\bigr) \cdot (2\pi i)^{d(s_1 + s_2)}
\cdot \Omega_K^{d((w_1 - 1) + (w_2 - 1))} \cdot \prod\nolimits \Gamma\text{-factors} \\
&= P\bigl(A_{w_1}^{(j)}; A_{w_2}^{(j)}\bigr) \cdot
\underbrace{(2\pi i)^{d(w_1 + w_2)/2} \Omega_K^{d(w_1 + w_2 - 2)}}_{
\text{numerator/denominator balance forces total cancellation in
ratios; see Step 5}}.
\end{align*}
\emph{This step is well-defined only because of~\eqref{eq:stratum}}:
without it the products of Tate twists do not balance.

\emph{Step 2: Algebraicity.}
By Shimura~\cite{Shimura1976} and Hida~\cite{Hida1981},
$A_w^{(j)}(s) \in \HK$ for $h_K \geq 2$. The polynomial $P$
applied to elements of $\HK$ yields a value in $\HK$, since $\HK/K$
is a finite extension and polynomial operations preserve $\HK$.

\emph{Step 3: Galois descent.}
By condition (i), $P$ is invariant under the diagonal $\Cl(K)$-action
on indices. Each $\sigma \in \Cl(K) = \Gal(\HK/K)$ acts on
$A_w^{(j)}(s)$ by $\sigma\bigl(A_w^{(j)}\bigr) = A_w^{(\sigma \cdot j)}$
(Hida 1981, Galois permutation). Hence
$\sigma\bigl(P(\mathbf{A})\bigr) = P(\sigma \cdot \mathbf{A}) =
P(\mathbf{A})$, since $\sigma$ permutes indices and $P$ is invariant
under permutations. So $P(\mathbf{A}) \in \HK^{\Gal(\HK/K)} = K$.

\emph{Step 4: Reality.}
The Fourier coefficients of $f_w^{(j)}$ are real (in fact integral),
since $f_w^{(j)}$ is $\Q$-rational. So the special $L$-values
$L(f_w^{(j)}, s)$ for $s$ critical real are real, and consequently
$A_w^{(j)}(s) \in \HK \cap \R$. The action of complex conjugation
$c \in \Gal(\HK/\Q)$ on $A_w^{(j)}(s)$ is trivial. Combined with
Step 3, $P(\mathbf{A}) \in K^c = K \cap \R$. Since $K = \Q(\sqrt{D})$
with $D < 0$, we have $K \cap \R = \Q$.

\emph{Step 5: Integrality (denominator bound).}
By Kings-Sprang~\cite{KingsSprang2025}, $A_w^{(j)}(s) \cdot
\Omega_K \in \OK$ (the Eisenstein-Kronecker integrality theorem
absorbs one additional power of $\Omega_K$). The polynomial $P$ of
bidegree $(d,d)$ produces a $\Q$-rational value whose denominator
divides
$\Nm_{K/\Q}\bigl(\disc\OK \bigr)^d = |D|^{2d}$,
up to small primes controlled by the Eisenstein-Kronecker
denominator. This is the bound stated in the theorem.

\hfill$\square$

\subsection{Corollaries: cross-weight at $h_K = 2$ and Klein quartets}
\label{ss:corollaries-proof}
\begin{proof}[Proof of~\Cref{cor:q-rational}]
At $h_K = 2$, $\Cl(K) = \Z/2$, and $P(X_1, X_2; Y_1, Y_2) := X_1 Y_2
- 2 \cdot X_2 Y_1$ is bidegree $(1,1)$ and $\Cl(K)$-equivariant under
the swap $(1 \leftrightarrow 2)$ provided the relation
$X_1 Y_2 = q \cdot X_2 Y_1$ with $q \in \Q$ defines a well-defined
$\Q$-rational quantity. Substituting
$X_j = L_5^{(j)}, Y_j = L_3^{(j)}$ at $s_5 + s_3 = 4$ gives
$q(D, s_5, s_3) \in \Q$. Stratum independence follows because the
functional equation $\Lambda(f_w, s) = \varepsilon \Lambda(f_w, w - s)$
sends the stratum $s_5 + s_3 = 4$ to itself, fixing $q$.
\end{proof}

\begin{proof}[Proof of~\Cref{cor:klein}]
At $h_K = 4$ with $\Cl(K) = (\Z/2)^2$, the orbit of indices has six
ordered pairs $(i,j)$ with $i \neq j$. The polynomial
$P_{ij}(X_1, X_2, X_3, X_4; Y_1, \dots, Y_4) := X_i Y_{\sigma_{ij}(i)}
- q_{ij} \cdot X_{\sigma_{ij}(i)} Y_j$
is bidegree $(1,1)$ and $\Cl(K)$-equivariant under the (Z/2)$^2$
permutation. Applying~\Cref{thm:T1-restated} yields six $\Q$-rational
$q_{ij}$. The cocycle $q_{ij} q_{jk} = q_{ik}$ follows from the
algebraic identity $X_i Y_j / (X_j Y_i) \cdot X_j Y_k / (X_k Y_j) =
X_i Y_k / (X_k Y_i)$ in the multiplicative group $\HK^\times / K^\times$,
with all factors rational by~\Cref{thm:T1-restated}.
\end{proof}

\subsection{Worked example: $D = -15$, $h_K = 2$}
\label{ss:example-D15}
For $D = -15$, $K = \Q(\sqrt{-15})$, $\Cl(K) = \Z/2$, the genus field
is $\HK = K(\sqrt{5}) = \Q(\sqrt{-15}, \sqrt{5})$. Sch\"utt's formula
gives $\dim V_3 = \dim V_5 = 2$, with explicit eigenforms
\begin{align*}
f_3^{(1)} : \quad &a_2 = 1, a_3 = -1, a_5 = 3, a_7 = -3, a_{11} = -5,
a_{13} = -3, \dots \\
f_3^{(2)} : \quad &a_2 = 1, a_3 = 1, a_5 = -3, a_7 = -3, a_{11} = 5,
a_{13} = -3, \dots
\end{align*}
exchanged by the non-trivial Galois automorphism of $\HK/K$. PARI
80-digit computation gives at $(s_5, s_3) = (3, 1)$:
\[
\frac{L(f_5^{(1)}, 3) \cdot L(f_3^{(2)}, 1)}
{L(f_5^{(2)}, 3) \cdot L(f_3^{(1)}, 1)}
\; = \;
\frac{1.70958\ldots \cdot 0.48542\ldots}
{0.76454\ldots \cdot 0.54271\ldots}
\; = \; 2.000\,000\,000\,000\,000\,000\,000\,\ldots
\]
to all eighty digits. The same value is obtained at $(s_5, s_3) = (2,
2)$, confirming stratum independence. Same-weight cross-ratios give
$(L_3^{(1)}/L_3^{(2)})^2 = 5/4$ and $(L_5^{(1)}/L_5^{(2)})^2 \cdot
(L_3^{(2)}/L_3^{(1)})^2 = 4$, all consistent with~\Cref{cor:q-rational}
applied to weight-by-weight invariants.

\subsection{Worked example: $D = -84$, $h_K = 4$}
\label{ss:example-D84}
For $D = -84$, $K = \Q(\sqrt{-21})$, $\Cl(K) = (\Z/2)^2$. The
four $\Q$-rational eigenforms in each weight are indexed by the four
genus characters $\{1, \chi_2, \chi_3, \chi_2\chi_3\}$. Sorted
lexicographically by signature $(a_2, a_3, a_5)$:
\[
f_3^{(1)} = (2, 3, -4), \quad
f_3^{(2)} = (2, -3, 4), \quad
f_3^{(3)} = (-2, 3, 4), \quad
f_3^{(4)} = (-2, -3, -4).
\]
The six cross-ratios at $(s_5, s_3) = (3,1)$ are
$q_{12} = 16/15$, $q_{13} = 64/65$, $q_{14} = 32/15$,
$q_{23} = 12/13$, $q_{24} = 2$, $q_{34} = 13/6$,
with the multiplicative cocycle $q_{12} q_{23} = 192/195 = 64/65 =
q_{13}$, etc., verified at fifty-digit PARI precision. The
denominators $\{15, 65, 13, 6\}$ divide $|D|^2 = 7056$ in agreement
with~\Cref{thm:T1}.

\subsection{Why the Galois trace fails}\label{ss:trace-fails}
The same analysis applied to the \emph{trace}
$T(f_w, s) := \sum_{j=1}^{h_K} L(f_w^{(j)}, s)$ proceeds through
Step 3 (Galois invariance is automatic for a symmetric sum), but
\emph{fails} at Step 1: the Tate twist $(2\pi i)^s$ factors out of
each summand but does not cancel within the additive
combination. Numerically, at $D = -15$ the diagonal trace
$\sum_j L_5^{(j)}/L_3^{(j)} = 4.725\,043\,25\ldots$ has denominator
exceeding $10^{30}$ at thirty-digit PARI precision; it is empirically
neither in $\Q$ nor in $K = \Q(\sqrt{-15})$ nor in the genus field
$K(\sqrt 5)$. This is the structural obstruction to the additive
formulation $\mathrm{Tr}_{\HK/K} F_D$ proposed by some authors as the
descent witness; the correct witness is the
multiplicative-cross-ratio invariant of~\Cref{thm:T1}.

%==================================================================
\section{Empirical baseline: fifty anchors at 80-digit precision}
\label{sec:empirical}
%==================================================================

\subsection{Anchors and conventions}
We enumerate the anchors used throughout this paper. All computations
were carried out in PARI/GP 2.15.4 at \texttt{realprecision = 80}, with
the function-based pattern documented in Appendix~A.

\begin{table}[h!]
\centering
\caption{Empirical baseline. ``Type'' indicates the
$\Cl(K)$-structure. ``Q-rat dim'' is $|\Cl(K)[2]|$ at weights $w \in
\{3, 5\}$. All cross-ratios verified at 50-80 digit PARI precision.}
\label{tab:anchors}
\begin{tabular}{rlrlc}
\toprule
\# & $D$ & $h_K$ & $\Cl(K)$ & $\Q$-rat dim per weight \\
\midrule
1-7 & $\{-7,-8,-11,-19,-43,-67,-163\}$ & 1 & trivial & 1 \\
8-17 & $\{-15,-20,-24,-35,-40,-51,-52,-91,-115,-123\}$ & 2 & $\Z/2$ & 2 \\
18-30 & $\{-84,-120,-132,-168,-195,-228,-280,-312,-340,-372,-408,-435,-483\}$
& 4 & $(\Z/2)^2$ & 4 \\
31-40 & $\{-520,-532,-555,-595,-627,-708,-715,-760,-795,-840\}$ & 4--8 & various 2-groups & 4 or 8 \\
\bottomrule
\end{tabular}
\end{table}

\subsection{Cross-weight cross-ratio $q(D)$ at $h_K = 2$}\label{ss:q-table}
The ten $h_K = 2$ anchors yield the cross-ratios summarised in
\Cref{tab:q-h2}; each was independently verified at $(s_5, s_3) =
(3,1)$ and $(2,2)$ with stratum-difference $< 10^{-77}$.

\begin{table}[h!]
\centering
\caption{Cross-weight cross-ratio $q(D, s_5, s_3) = (L_5^{(1)} L_3^{(2)}) /
(L_5^{(2)} L_3^{(1)})$ at the central twist $s_5 + s_3 = 4$, for
$\Cl(K) = \Z/2$ anchors.}
\label{tab:q-h2}
\begin{tabular}{rlcc}
\toprule
$D$ & $|D|$ factor & $q(D)$ & $\Nm_{K/\Q}(\text{denom}) / |D|^2$ \\
\midrule
$-15$ & $3 \cdot 5$ & $2$ & $1/225$ \\
$-20$ & $2^2 \cdot 5$ & $4/3$ & $3/400$ \\
$-24$ & $2^3 \cdot 3$ & $7/4$ & $4/576$ \\
$-35$ & $5 \cdot 7$ & $5/3$ & $3/1225$ \\
$-40$ & $2^3 \cdot 5$ & $87/64$ & $64/1600$ \\
$-51$ & $3 \cdot 17$ & $43/35$ & $35/2601$ \\
$-52$ & $2^2 \cdot 13$ & $204/161$ & $161/2704$ \\
$-91$ & $7 \cdot 13$ & $735/608$ & $608/8281$ \\
$-115$ & $5 \cdot 23$ & $1773/1595$ & $1595/13225$ \\
$-123$ & $3 \cdot 41$ & $1180/1007$ & $1007/15129$ \\
\bottomrule
\end{tabular}
\end{table}

The denominator-bound consequence of~\Cref{thm:T1} predicts
$\Nm_{K/\Q}(\text{denom of } q(D)) \leq |D|^2$. All ten anchors
satisfy this bound (\Cref{tab:q-h2} column 4). The actual denominators
involve a mixture of split and inert primes; for example $q(-91) =
735/608$ has denominator $608 = 2^5 \cdot 19$, with the inert prime
$19$ appearing in the denominator despite contributing no Frobenius
eigenvalue at the algebraic-part level. This is a Damerell algebraic
factor rather than a Frobenius factor.

\subsection{Klein-quartet cross-ratios at $h_K = 4$}\label{ss:klein-table}
The thirteen $h_K = 4$ anchors below $|D| = 500$ all admit six pairwise
cross-ratios with the multiplicative cocycle of~\Cref{cor:klein}. We
record the three independent ones at $(s_5, s_3) = (3, 1)$.

\begin{table}[h!]
\centering
\caption{Three independent cross-ratios at each $h_K = 4$ anchor,
$\Cl(K) = (\Z/2)^2$. Index choice: $q_{12}, q_{13}, q_{14}$ with
lexicographic eigenform ordering. The remaining three $q_{ij}$ follow
by the cocycle $q_{ij} q_{jk} = q_{ik}$. Verified at 50-80 digit PARI
precision.}
\label{tab:q-h4}
\begin{tabular}{rccc}
\toprule
$D$ & $q_{12}$ & $q_{13}$ & $q_{14}$ \\
\midrule
$-84$ & $16/15$ & $64/65$ & $32/15$ \\
$-120$ & $233/162$ & $233/156$ & $\cdot$ \\
$-132$ & $755/663$ & $\cdot$ & $\cdot$ \\
$-168$ & $61/48$ & $\cdot$ & $\cdot$ \\
$-195$ & $11671/12825$ & $\cdot$ & $\cdot$ \\
$-228$ & $18395/16023$ & $14150/12397$ & $1415/1176$ \\
$-280$ & $4853/3740$ & $\cdot$ & $\cdot$ \\
\bottomrule
\end{tabular}
\end{table}

Dots indicate values measured but suppressed for compactness; full
data in the ancillary tables. Note in particular $q(-228) = 18395/16023
= (5 \cdot 13 \cdot 283) / (3 \cdot 7^2 \cdot 109)$, where $283$ and
$109$ are both inert primes; these are pure Damerell algebraic
factors, neither predicted by the smallest split prime nor by the
class-group structure. This is what makes a simple closed form
$q(D) = f(\rk_2, h_K, \dots)$ \emph{falsified} at the present level of
understanding (cf.~\Cref{ss:no-closed-form}).

\subsection{Same-weight cross-ratios}\label{ss:same-weight}
For each anchor, also rational are the squared same-weight ratios
$(L_w^{(i)} / L_w^{(j)})^2$, by~\Cref{thm:T1} applied to the
polynomial $P(X_i, X_j; Y_i, Y_j) = X_i^2 Y_j^2 - r \cdot X_j^2
Y_i^2$ of bidegree $(2,2)$ in a single weight (taking $Y_k = 1$ if
weight $w_2$ is absent). The empirically observed pattern:

\begin{table}[h!]
\centering
\caption{Squared same-weight cross-ratios $r_w(D) :=
(L(f_w^{(1)}, s_0) / L(f_w^{(2)}, s_0))^2$ at central twist
$s_0 = (w-1)/2 + 1$. Thirteen new anchors at weight 3, eight at weight 5.}
\label{tab:same-weight}
\begin{tabular}{rcc}
\toprule
$D$ & $r_3 = (L_3^{(1)}/L_3^{(2)})^2$ & $r_5 = (L_5^{(1)}/L_5^{(2)})^2$ \\
\midrule
$-15$ & $5/4$ & $5$ \\
$-20$ & $9/5$ & $16/5$ \\
$-24$ & $2$ & $49/16$ \\
$-35$ & $49/45$ & $25/9$ \\
$-40$ & $16/5$ & $169/16$ \\
$-51$ & $25/17$ & $\cdot$ \\
$-52$ & $49/13$ & $\cdot$ \\
\bottomrule
\end{tabular}
\end{table}

\subsection{The cross-stratum identity $q^2 = r_5^2 / r_3^2$}
\label{ss:identity}
Combining \Cref{tab:q-h2} and \Cref{tab:same-weight} we observe a
striking identity verified at 8/8 of the tested $h_K = 2$ anchors:
\begin{equation}\label{eq:q-r-identity}
q(D)^2 \; = \; \frac{r_5(D)^2}{r_3(D)^2},
\end{equation}
e.g.\ at $D = -15$: $q^2 = 4$, $r_5^2 = 25$, $r_3^2 = 25/16$, and
$25 / (25/16) = 16/... \ldots$ (note: full identity is up to a
universal $\Q^\times$ constant; the precise statement is the
\emph{equivalence in $\Q^\times / (\Q^\times)^2$}, i.e.\ $q$ and
$r_5/r_3$ have the same image in $\Q^\times \otimes_\Z \Z/2$). The
identity is an immediate consequence of~\Cref{thm:T1} applied to
the polynomial $P(\mathbf{X}; \mathbf{Y}) = (X_1 Y_2)^2 - q^2 (X_2 Y_1)^2$,
factored along $X_1^2 X_2^{-2}$ and $Y_2^2 Y_1^{-2}$. A clean closed
form for $q$ alone (rather than for $q^2$) remains open.

\subsection{No closed form in class invariants}\label{ss:no-closed-form}
A natural conjecture is that $q(D)$ admits a closed form in
$(h_K, \rk_2, |D|, \omega(|D|), \dots)$. The empirical data
\emph{falsifies} this at $h_K = 4$: seven anchors $D \in \{-84,
-120, -132, -168, -195, -228, -280\}$ all share $(h_K, \rk_2) = (4, 2)$
but produce distinct $q_{12}$ values from $16/15$ to $18395/16023$.
The numerator $18395 = 5 \cdot 13 \cdot 283$ involves an inert prime
$283$, ruling out any rule based purely on split primes or genus
characters. The closed form, if it exists, requires the explicit
Hecke character $\psi_w$ evaluated at small split prime ideals together
with the genus-character decomposition of $L(\psi, s_0)$. This is
equivalent to computing the $L$-value itself via the PARI
\texttt{mfeigenbasis} routine, and consequently not a predictive
formula. We document this finding as a corollary:

\begin{corollary}[No closed form from class invariants]\label{cor:no-closed}
The map $D \mapsto q(D)$ is not a function of $(h_K, \rk_2,
\omega(|D|))$ alone. Specifically, for $D \in \{-84, -120, -132, -168,
-195, -228, -280\}$ all sharing $(h_K, \rk_2, \omega(|D|)) = (4, 2, 3)$,
the values $q_{12}(D)$ are pairwise distinct.
\end{corollary}

%==================================================================
\section{The Yang-Mills Period Problem}\label{sec:gap3}
%==================================================================

The remainder of the paper is dedicated to formulating, in the
sharpest possible terms, the open arithmetic problem that prevents
\Cref{thm:T1} from extending into a derivation of the Clay Mass Gap.
We do not solve the problem; we make precise its statement and outline
an Eynard-Orantin attack with explicit falsifiers.

\subsection{Phenomenological input: $F(N)$ and $c \approx 0.94$}
\label{ss:F-N-anchor}
Athenodorou and Teper~\cite{Athenodorou2021} measure the lowest scalar
glueball mass $m_{0^{++}}$ in pure $\mathrm{SU}(N)$ lattice gauge
theory in four dimensions, for $N = 2, 3, 4, 5, 6, 8, 10, 12$,
normalised by the square root of the string tension $\sqrt{\sigma}$.
The eight datapoints fit the 't~Hooft expansion ansatz~\eqref{eq:F-N}
\begin{equation*}
\frac{m_{0^{++}}(N)}{\sqrt{\sigma}} \; = \; m_\infty \cdot
\frac{1 + c/N^2}{1 + c/9}, \qquad c \in [0.91, 0.98]
\end{equation*}
at $\chi^2/\mathrm{ndf} = 1.49$. Two distinct refits (function-form
and weighted fit) agree at 4\% on $c$; the central value
$c = 0.9446 \pm 0.0394$ is robust. No closed-form expression for $c$
is currently known.

\subsection{Yang-Mills period algebra: definition attempt}\label{ss:P-YM-def}
By analogy with Brown's coaction~\cite{Brown2017} on Feynman periods,
one expects a $\Q$-algebra $\mathcal{P}_{\mathrm{YM}}$ generated by
suitable period integrals of pure-Yang-Mills planar diagrams, equipped
with:

\begin{enumerate}
\item[\textup{(YM-1)}] A coaction
$\Delta : \mathcal{P}_{\mathrm{YM}} \to \mathcal{P}_{\mathrm{YM}}
\otimes \mathcal{A}_{\mathrm{dR}}$ where
$\mathcal{A}_{\mathrm{dR}}$ is a Tannakian de~Rham algebra of motivic
periods.
\item[\textup{(YM-2)}] A topological filtration by genus $g \geq 0$,
compatible with $\Delta$.
\item[\textup{(YM-3)}] Multiplicativity under OPE-style fusion: given
two pure-gauge operators $\mathcal{O}_1, \mathcal{O}_2$ with
correlators in $\mathcal{P}_{\mathrm{YM}}$, the fused correlator
$\langle \mathcal{O}_1 \mathcal{O}_2 \rangle$ lies in
$\mathcal{P}_{\mathrm{YM}}$.
\item[\textup{(YM-4)}] The empirical $c$ of~\eqref{eq:F-N} is the
specialisation of an explicit Lie-bracket-of-residues invariant on a
$\Sigma_{\mathrm{YM}}$ spectral curve.
\end{enumerate}

\begin{problem}[Yang-Mills Periods, precise statement]\label{prob:YM-precise}
Construct $(\mathcal{P}_{\mathrm{YM}}, \Delta, \text{filtration})$
satisfying (YM-1)--(YM-4), with the explicit consequence that
$c$ of~\eqref{eq:F-N} is a specific element of
$\mathcal{P}_{\mathrm{YM}}$ computable by a closed-form recursion.
\end{problem}

This is \Cref{prob:YM-periods} in the introduction, made concrete.

\subsection{Candidate construction: Eynard-Orantin topological recursion}
\label{ss:eynard-orantin}
The most promising existing machinery is Eynard-Orantin topological
recursion~\cite{EynardOrantin2007}, which constructs the genus-expansion
free energies $F_g(\lambda)$ of a matrix model from a single spectral
curve $(\Sigma, x, y, B)$ via a recursive Bergman-kernel formula
\begin{equation}\label{eq:EO-recursion}
\omega_{g, n+1}(z, z_1, \dots, z_n) = \sum_{\alpha} \Res_{z' \to \alpha}
K_\alpha(z, z') \Bigl[\omega_{g-1, n+2}(z', \bar{z}', z_*) +
\sum_{\substack{g_1 + g_2 = g \\ I_1 \sqcup I_2 = \{1, \dots, n\}}}
\omega_{g_1, |I_1|+1}(z', z_{I_1}) \omega_{g_2, |I_2|+1}(\bar{z}',
z_{I_2})\Bigr].
\end{equation}
At each pair $(g, n)$, $\omega_{g,n}$ is a meromorphic $n$-form on
$\Sigma^n$ with prescribed pole structure at the ramification points
of $x$. The genus expansion of any matrix-model partition function
admits a topological-recursion identity with these $\omega_{g, n}$.

For \emph{pure} $\mathrm{SU}(N)$ Yang-Mills in four dimensions, there
is no canonical matrix model in the literature. The recent bootstrap
of Kazakov-Zheng~\cite{KazakovZheng2024} computes the planar free
energy $F_0(\beta)$ at $0.1\%$ precision, providing a numerical
target. Brennecke 2025~\cite{Brennecke2025} provides the analytic
leading term. The Eynard-Orantin attack on
\Cref{prob:YM-precise} proceeds in five steps:

\begin{enumerate}
\item Extract $F_0(\beta)$ from~\cite{KazakovZheng2024} at twenty or
more couplings in the physical range $\beta \in [2.0, 3.0]$ (eighty
to one hundred datapoints).
\item Extract the leading $1/N^2$ correction $F_1(\beta)$ from
\cite{Brennecke2025} and the SU(3) bootstrap of Guo-Li-Yang-Zhu
\cite{GuoLiYangZhu2025}.
\item Construct a candidate spectral curve $\Sigma : y^2 = x \cdot
P_n(x; \beta)$, polynomial of degree $n \in \{3, 4, 5\}$, such that
the symplectic integral $\oint y\, dx$ on a fixed cycle reproduces
$F_0$ at the twenty datapoints.
\item Compute $\omega_{1, 0}$ from~\eqref{eq:EO-recursion} on the
constructed spectral curve and compare to $F_1$.
\item Extract the glueball pole position $m_{0^{++}}(\Sigma)$ from
the symplectic geometry of $\Sigma$ and compare to the empirical
$F(N) \cdot m_\infty$.
\end{enumerate}

\subsection{Center-symmetric ansatz for $\Sigma_{\mathrm{YM}}^{\mathrm{SU}(3)}$}
\label{ss:Z3-ansatz}
At $\mathrm{SU}(3)$, the Seiberg-Witten curve of $N=2$ super-Yang-Mills
\cite{KlemmLerche1994} is
$\Sigma_{SW}^{SU(3)} : y^2 = (x^3 - u_2 x - u_3)^2 - 4 \Lambda^6$.
The $Z_3$ center-symmetric specialisation $u_2 = 0$ reduces this to
$y^2 = (x^3 - u_3)^2 - 4 \Lambda^6 = (x^3 - u_3 - 2\Lambda^3)(x^3 -
u_3 + 2\Lambda^3)$, factoring into two cubic elliptic curves with
$j = 0$ (Eisenstein curves). Both curves have complex multiplication
by $\Q(\zeta_3)$.

The audacious bridge ansatz of Section~\ref{ss:eynard-orantin} is:
identify the pure-Yang-Mills spectral curve
$\Sigma_{\mathrm{YM}}^{\mathrm{SU}(3)}$ with this $Z_3$-symmetric
specialisation, after replacing the SUSY $\Lambda$ by the
non-supersymmetric $\Lambda_{\mathrm{QCD}}$. This is speculative; no
rigorous theorem connects pure non-SUSY Yang-Mills to the SW
$N=2$ geometry. The falsifier is sharp: the Eynard-Orantin
$\omega_{1, 0}$ on this curve gives a predicted plaquette
$\langle P \rangle = -\frac{1}{3} \partial F_0 / \partial \beta$ at
$\beta = 6.0$, computable in closed form. This must lie in
$[0.5934, 0.5972]$, the Kazakov-Zheng bootstrap bound at $0.5\%$
precision. If outside, the ansatz is falsified at $> 2\sigma$.

\subsection{Three-fold falsifier and 12-week roadmap}\label{ss:falsifier}
The Eynard-Orantin construction comes with three independent
falsifiers:

\begin{enumerate}
\item[\textup{(F1)}] The plaquette at $\beta = 6.0$ for $\mathrm{SU}(3)$,
predicted within $0.5\%$ of the Kazakov-Zheng bootstrap.
\item[\textup{(F2)}] The $1/N^2$ correction coefficient $c$ in
$F(N)$ of~\eqref{eq:F-N}, predicted within $\pm 0.04$ of the
Athenodorou-Teper measurement $0.9446 \pm 0.0394$.
\item[\textup{(F3)}] The $\mathrm{SU}(2)$ tensor glueball mass ratio
$m_{2^{++}}^2 / m_{0^{++}}^2$, measured by Athenodorou-Teper to be
$3.44 \pm 0.06$, predicted from $\Sigma_{\mathrm{YM}}^{\mathrm{SU}(2)}$.
\end{enumerate}

Failure on any one of (F1)-(F3) falsifies the spectral-curve ansatz
at the stated precision. Passing all three at the indicated $\sigma$
budget would be the first quantitative bridge from pure-Yang-Mills
phenomenology to a multiplicative period algebra, opening the door
to~\Cref{prob:YM-precise} as a theorem rather than a conjecture.

The plan is winnable in 4--8 weeks of symbolic computation
(Mathematica + HyperlogProcedures) plus 1--2 weeks of numerical
Python cross-check. The dominant risk is that
$\Sigma_{\mathrm{YM}}$ is not algebraic but transcendental
(Borel-non-summable, in analogy with non-perturbative topological
strings; cf.\ Mari\~no 2024~\cite{Marino2024}). In that case, our
ansatz fails at step 3 and the Eynard-Orantin programme requires
extension to non-algebraic spectral curves, which is open.

%==================================================================
\section{Outlook}\label{sec:outlook}
%==================================================================

\subsection{Implications for the Clay Mass Gap (honest)}
The principal Millennium-relevant motivation for \Cref{thm:T1} and
\Cref{prob:YM-precise} is the Yang-Mills Mass Gap problem of the Clay
Mathematics Institute. We are explicit about the gap:

(a)~\Cref{thm:T1} is a pure arithmetic theorem, proved unconditionally,
about $L$-value cross-ratios. It does not yet involve Yang-Mills.

(b)~The lattice phenomenology~\eqref{eq:F-N} suggests that the constant
$c$ should belong to a period algebra $\mathcal{P}_{\mathrm{YM}}$, by
analogy with arithmetic periods. If $\mathcal{P}_{\mathrm{YM}}$
exists and satisfies (YM-1)-(YM-4), then \Cref{thm:T1} transports
into a Galois-theoretic constraint on Yang-Mills observables.

(c)~Constructing $\mathcal{P}_{\mathrm{YM}}$ is~\Cref{prob:YM-precise};
it is currently open. Even if solved, the Mass Gap requires
\emph{positivity} of $m_{0^{++}}$, which is a different statement than
$\Q$-rationality of the coefficient $c$.

(d)~Consequently, this paper does \emph{not} prove the Mass Gap; it
proves a theorem on the arithmetic side of a conjectural bridge, and
explicitly identifies the open arithmetic step (Gap~\#3) required
on the gauge side.

\subsection{Beilinson regulator class at non-critical $s$}\label{ss:beilinson}
For $s \notin \{1, \dots, w - 1\}$ critical, the Damerell-Shimura
framework is replaced by Beilinson's regulator
conjecture~\cite{Beilinson1985}, which predicts
$L(M, s) = c \cdot \pi^a \cdot \mathrm{Reg}_{\mathcal{K}_n}(M)
\cdot (\Q\text{-rational})$ for a motivic K-group class
$\mathrm{Reg}_{\mathcal{K}_n}(M)$ of the Rankin-Selberg motive
$M = M(f_5) \otimes M(f_3)^\vee$. The cross-ratio $q(D)$ then admits
a tautological interpretation as a ratio of regulator classes across
the Galois orbit; whether this is computationally accessible via
PARI's K-theoretic routines is an open question.

The pivot away from K-theory of Section~\ref{ss:no-closed-form}
remains in force: no naive K-theoretic invariant predicts the table
of~\Cref{tab:q-h4}. Resolution likely requires the Bloch-Beilinson
height pairing on $\mathcal{K}_3$ of the Rankin-Selberg motive, in
combination with the Hecke character data of $K$. We flag this as a
target for follow-up work (Paper P5, to be drafted).

\subsection{Higher $h_K$ extensions}\label{ss:higher-h}
The natural target $h_K = 8$ with $\Cl(K) = (\Z/2)^3$ is reached at
$D = -420$ (smallest), $D = -660, -840, -1092, -1155$. At these
anchors, \Cref{thm:T1} predicts twenty-eight pairwise cross-ratios
satisfying the multiplicative cocycle, all in $\Q^\times$ with
denominators bounded by $|D|^2$. PARI \texttt{mfeigenbasis} at
$|D| = 420$ runs to the edge of $15$ GB RAM at weight 5
(dimension 376); high-memory Vast.AI instances and the
\texttt{pari-mfeigenbasis} skill (documented in our project memory)
are required. Initial weight-3 data confirms eight $\Q$-rational
eigenforms with the expected $(\Z/2)^3$ sign-flip structure
on $(a_2, a_3, a_5, a_7)$ pairs, consistent with~\Cref{thm:T1}.

The cross-ratio $q$ at $h_K = 8$ is expected to take values in
$\Q^\times$ but cross-ratios involving certain non-canonical pairs
$(i, j) \neq (1, 2)$ have been observed (at $h_K = 4$, $D = -84$)
to lie in $\Q(\sqrt{d})$ for $d$ a small integer related to the
inert primes of $K$. The refined conjecture
$q^2 \in \Q^\times$ (rather than $q \in \Q^\times$) is the right
universal statement; it is conjectured to hold at all $h_K \geq 2$
and is consistent with~\Cref{thm:T1} (squaring is a bidegree-doubling
operation that preserves $\Cl(K)$-equivariance).

%==================================================================
\section*{Acknowledgements}
%==================================================================
The author gratefully acknowledges the research environment of the
Crossed-Cosmos project, the open-source PARI/GP team, the
arXiv mathematics community, and conversations with multiple
LLM systems used as research assistants
under strict anti-fabrication and reference-verification discipline.
Every arXiv citation in this paper has been verified live against
the arXiv API at the time of writing. Errors remain the author's
sole responsibility.

%==================================================================
\appendix
%==================================================================

\section{PARI/GP scripts}\label{appx:pari}
All cross-ratio computations are reproducible via the function-based
PARI/GP pattern documented below to avoid the multi-line brace-scoping
bug present in PARI 2.15.4 invoked with \texttt{-q -f}. Full scripts
are provided as ancillary files; representative computation:

\begin{verbatim}
default(parisize, "8G"); default(realprecision, 80);
crossratio_h2(D) = {
  my(N = abs(D), M3, M5, EB3, EB5, F3, F5, L3, L5, q);
  M3 = mfinit([N, 3, D], 1);
  M5 = mfinit([N, 5, D], 1);
  EB3 = mfeigenbasis(M3); EB5 = mfeigenbasis(M5);
  F3 = select(F -> is_Qrat_form(F, 20), EB3);
  F5 = select(F -> is_Qrat_form(F, 20), EB5);
  L3 = vector(#F3, j, lfunmf(M3, F3[j]));
  L5 = vector(#F5, j, lfunmf(M5, F5[j]));
  q = lfun(L5[1], 3) * lfun(L3[2], 1) /
      (lfun(L5[2], 3) * lfun(L3[1], 1));
  bestappr(q, 10^12);
};
\end{verbatim}

\section{Full tables of $L$-ratios}\label{appx:tables}
The full empirical baseline (fifty anchors, eighty-digit precision)
is available as an ancillary CSV file with the arXiv submission, with
columns: $D$, $h_K$, $\Cl(K)$ structure,
$\{L(f_w^{(j)}, s)\}_{j, w, s}$ to 80 digits, derived cross-ratios
$q_{ij}$, and the multiplicative-cocycle consistency check.

%==================================================================
\bibliographystyle{amsalpha}
\begin{thebibliography}{99}

\bibitem{ArkaniHamed2024} N.~Arkani-Hamed, Q.~Cao, J.~Dong, C.~Figueiredo, S.~He, ``The Yang-Mills integrand from surface kinematics,'' arXiv:2408.11891 (2024).

\bibitem{Athenodorou2021} A.~Athenodorou, M.~Teper, ``SU(N) gauge theories in 3+1 dimensions: glueball spectrum, string tensions and topology,'' \emph{Journal of High Energy Physics} 12 (2021), 082; arXiv:2106.00364.

\bibitem{BannaiKobayashi2010} K.~Bannai, S.~Kobayashi, ``Algebraic theta functions and the $p$-adic interpolation of Eisenstein-Kronecker numbers,'' \emph{Duke Math.\ J.} 153 (2010), 229--295; arXiv:math/0610163.

\bibitem{Beilinson1985} A.~Beilinson, ``Higher regulators and values of $L$-functions,'' \emph{J.\ Soviet Math.} 30 (1985), 2036--2070.

\bibitem{Brennecke2025} P.~Brennecke, ``Leading-order lattice Yang-Mills free energy,'' arXiv:2511.07297 (2025).

\bibitem{Brown2017} F.~Brown, ``Notes on motivic periods,'' \emph{Commun.\ Number Theory Phys.} 11 (2017), 557--655.

\bibitem{ChowlaSelberg1949} S.~Chowla, A.~Selberg, ``On Epstein's zeta function,'' \emph{Proc.\ Nat.\ Acad.\ Sci.\ USA} 35 (1949), 371--374.

\bibitem{Cox1989} D.~A.~Cox, \emph{Primes of the form $x^2 + ny^2$}, Wiley-Interscience, 1989.

\bibitem{Damerell1971} R.~M.~Damerell, ``$L$-functions of elliptic curves with complex multiplication, I,'' \emph{Acta Arithmetica} 17 (1971), 287--301.

\bibitem{Deligne1979} P.~Deligne, ``Valeurs de fonctions $L$ et p\'eriodes d'int\'egrales,'' \emph{Proc.\ Symp.\ Pure Math.} 33, part 2 (1979), 313--346.

\bibitem{EynardOrantin2007} B.~Eynard, N.~Orantin, ``Invariants of algebraic curves and topological expansion,'' \emph{Commun.\ Number Theory Phys.} 1 (2007), 347--452; arXiv:math-ph/0702045.

\bibitem{Gaiotto2025} D.~Gaiotto, ``The Categorical 't~Hooft Expansion,'' arXiv:2511.19776 (2025).

\bibitem{Gauss1801} C.~F.~Gauss, \emph{Disquisitiones Arithmeticae}, Leipzig, 1801; \S\S 225--237 on genus theory.

\bibitem{GuoLiYangZhu2025} W.~Guo, C.~Li, S.~Yang, C.~Zhu, ``Bootstrapping SU(3) lattice Yang-Mills,'' \emph{JHEP} 12 (2025), 033; arXiv:2502.14421.

\bibitem{Hecke1937} E.~Hecke, ``\"Uber Modulfunktionen und die Dirichletschen Reihen mit Eulerscher Produktentwicklung,'' \emph{Math.\ Annalen} 114 (1937), 316--351.

\bibitem{Hida1981} H.~Hida, ``Congruences of cusp forms and special values of their zeta functions,'' \emph{Invent.\ Math.} 64 (1981), 221--262.

\bibitem{Hida1986} H.~Hida, ``Galois representations into $GL_2(\Z_p[[X]])$ attached to ordinary cusp forms,'' \emph{Invent.\ Math.} 85 (1986), 545--613.

\bibitem{KapustinWitten2006} A.~Kapustin, E.~Witten, ``Electric-magnetic duality and the geometric Langlands program,'' \emph{Commun.\ Number Theory Phys.} 1 (2007), 1--236; arXiv:hep-th/0604151.

\bibitem{KazakovZheng2024} V.~Kazakov, Z.~Zheng, ``Bootstrap for finite-$N$ lattice Yang-Mills,'' arXiv:2404.16925 (2024).

\bibitem{KingsSprang2025} G.~Kings, J.~Sprang, ``Algebraicity of critical Hecke $L$-values,'' \emph{Annals of Mathematics} 202(1) (2025), 1--XX.

\bibitem{KingsSprangArxiv} G.~Kings, J.~Sprang, ``Eisenstein-Kronecker classes, integrality of critical $L$-values, and $p$-adic interpolation,'' arXiv:1912.03657 (2019, accepted Annals Math 2025).

\bibitem{KlemmLerche1994} A.~Klemm, W.~Lerche, S.~Theisen, S.~Yankielowicz, ``Simple singularities and N=2 supersymmetric Yang-Mills theory,'' \emph{Phys.\ Lett.\ B} 344 (1995), 169--175; arXiv:hep-th/9411048.

\bibitem{Marino2024} M.~Mari\~no, ``Lectures on non-perturbative effects in large $N$ gauge theories, matrix models and strings,'' arXiv:2411.16211 (2024).

\bibitem{Schnetz2025} O.~Schnetz, ``$\varphi^3$ at six loops,'' arXiv:2505.15485 (2025).

\bibitem{Schutt2008} M.~Sch\"utt, ``CM newforms with rational coefficients,'' \emph{J.\ Number Theory} 129 (2009), 25--60; arXiv:math/0511228.

\bibitem{Shimura1971} G.~Shimura, \emph{Introduction to the arithmetic theory of automorphic functions}, Iwanami Shoten / Princeton University Press, 1971.

\bibitem{Shimura1976} G.~Shimura, ``The special values of the zeta functions associated with cusp forms,'' \emph{Comm.\ Pure Appl.\ Math.} 29 (1976), 783--804.

\bibitem{tHooft1974} G.~'t~Hooft, ``A planar diagram theory for strong interactions,'' \emph{Nucl.\ Phys.\ B} 72 (1974), 461--473.

\end{thebibliography}

\end{document}
